% =====================================================================
% Glosario. Redactado el 06/11/2026 (sesión R4), a mano como lista de
% descripciones (sin el paquete glossaries). Las definiciones son las del
% uso que se hace de cada término en la memoria.
% =====================================================================
\chapter*{Glosario}
\addcontentsline{toc}{chapter}{Glosario}

Este glosario recoge, por orden alfabético, los términos técnicos que más se
usan en la memoria, con el sentido que tienen en ella.

\begin{description}
	\item[ACI (inferencia conforme adaptativa).] Método que construye el intervalo
	de predicción con los residuos recientes del modelo y ajusta en línea su nivel
	según los fallos pasados, con un paso de adaptación $\gamma$.
	\item[ADF y KPSS.] Contrastes de estacionariedad con hipótesis nulas opuestas:
	raíz unitaria el ADF, estacionariedad el KPSS.
	\item[ARMA.] Modelo de los errores que combina autorregresión y media móvil. En
	el SARIMAX de este trabajo modela la dependencia de los errores de la regresión.
	\item[Calibración.] Coincidencia entre la frecuencia con que la observación cae
	dentro de un intervalo y el nivel nominal de este.
	\item[Cobertura empírica.] Proporción de observaciones que caen dentro del
	intervalo de predicción.
	\item[CRPS.] Puntuación de una distribución predictiva completa; aquí se aproxima
	con una rejilla de 19 cuantiles (CRPS$_{19}$).
	\item[Diebold--Mariano.] Contraste de igualdad de precisión entre dos
	previsiones, basado en la diferencia de sus pérdidas. Se usa con la corrección de
	Harvey, Leybourne y Newbold, la varianza de Newey y West y el ajuste de Holm.
	\item[ESIOS.] Sistema de información del operador del sistema eléctrico español
	(Red Eléctrica de España), de donde procede la demanda.
	\item[Fuga de información.] Uso, al ajustar o evaluar un modelo, de datos que no
	estarían disponibles en el momento de emitir la previsión.
	\item[Grados-día.] Medida de cuánto se aleja la temperatura de un umbral de
	confort, por debajo (calefacción) o por encima (refrigeración).
	\item[Huella de contenido.] Número calculado a partir de todos los valores del
	conjunto de datos, que cambia si cambia cualquiera de ellos.
	\item[Ingenuo estacional.] Previsión que copia la demanda de la misma hora del
	mismo día de la semana anterior.
	\item[LightGBM.] Implementación de \textit{gradient boosting}, que combina
	árboles de decisión añadidos de forma secuencial.
	\item[MAE, RMSE, MAPE y MASE.] Error absoluto medio, raíz del error cuadrático
	medio, error porcentual absoluto medio y error absoluto medio escalado. El MASE
	de este trabajo divide el MAE por el del ingenuo estacional en 2017--2022.
	\item[MSTL.] Descomposición de una serie en tendencia, varias estacionalidades y
	un resto.
	\item[N-HiTS.] Red neuronal de previsión que genera la previsión a varias
	escalas de tiempo.
	\item[Nitidez.] Estrechez de los intervalos de predicción: a igual calibración,
	son mejores los más estrechos.
	\item[Origen móvil.] Esquema de evaluación en el que la previsión se emite cada
	día con la información disponible hasta ese momento.
	\item[Periodograma.] Estimación de la potencia de una serie en cada frecuencia,
	que revela sus periodicidades.
	\item[Pérdida de pinball.] Pérdida asimétrica con la que se estima un cuantil.
	\item[Previsión meteorológica perfecta.] Supuesto de que la temperatura usada en
	la previsión es la observada, sin error de previsión.
	\item[Puntuación de intervalo.] Suma de la anchura de un intervalo y una
	penalización proporcional a lo que la observación queda fuera de él.
	\item[Reajuste.] Nueva estimación de un modelo con los datos disponibles hasta una
	fecha. En 2024 cada modelo se reajusta cada mes.
	\item[Reanálisis.] Reconstrucción de las condiciones meteorológicas pasadas
	combinando observaciones y un modelo físico.
	\item[SARIMAX.] En este trabajo, una regresión con términos de Fourier, variables
	exógenas y errores ARMA, estimada con la clase SARIMAX de statsmodels.
	\item[Sesgo.] Media de la demanda real menos la prevista; un valor negativo
	indica que se prevé más demanda de la que hubo.
	\item[Términos de Fourier.] Senos y cosenos con un periodo dado, que representan
	una estacionalidad en una regresión.
\end{description}
